\ssscp{Transcription in \citet{Onvlee1984}}{03-03-02}
One of our principal sources of data for several languages of
Sumba is \citet{Onvlee1984}, a dictionary focussed on Kambera.\footnote{
		The data for \citet{Onvlee1984}
		was collected between 1920--1950 \citep[259]{Klamer2009}.}
While Kambera
is the focus of \citeauthor{Onvlee1984}’s description,
cognates and/or semantically equivalent forms from six other
languages/dialects of Sumba are given in the etymological notes:
Mangili, Lewa, Anakalang, Mamboru, Weyewa, and Kodi.
While this dictionary is an invaluable source of data,
it has transcriptional issues which are discussed by \citet{Klamer2009}.

The main issue in the transcription of \citet{Onvlee1984} concerns the vowels.
The transcription of vowels is not phonemic
and mixes both phonetic and phonemic representation.
For instance, the acute accent is used to indicate a phonetically long vowel; e.g.
‹ú› = [uː] and ‹í› = [iː].
In some cases such length is contrastive,
while in other cases it is automatically predictable.
To illustrate, the final long vowel in ‹paní› ‘big bat’ (p.
398) is analysable as a final double vowel; /panii/ (from *paniki
with *k > {\0}).
This contrasts with single vowels in forms such as ‹pàni›
[pɐni] ‘chicken feed’ (from *paniŋ ‘bait, fodder’).
On the other hand,
the penultimate vowel in ‹tílu› ‘egg’ (p. 494) is also long,
but in this case the length is an automatic result of this vowel
bearing stress, as stressed vowels have increased length in Kambera \citep[252]{Klamer2009}.
Hence, the word ‘egg’ is analysable as /tilu/ (e.g.
see the examples transcribed ‹tilu› ‘egg’ in \citealp[148, 207]{Klamer1998a}).
As a result of such issues “[{\ldots}] \citet{Onvlee1984} is
not a reliable source for research on Kambera vowel length or
long vowel phonotactics, diachronic or synchronic, because the vowels are not represented
with a consistent and phonemic orthography.” \citep[253]{Klamer2009}.

Another issue in the transcription of \citet{Onvlee1984} concerns
the representation of lengthened consonants in Kodi and Weyewa.
\citeauthor{Onvlee1984} marks lengthened consonants with a grave accent over the
vowel preceding the consonant.
Thus, Kodi ‹lòdo› for /loɗːo/ ‘day’ (p.
243) and ‹lodo› for /loɗo/ ‘sing’ (p.246).\footnote{
		Recall that vowels preceding a lengthened consonant in Kodi are shorter
		than those preceding a single consonant; e.g.
		/loɗo/ [ˈlɔˑɗɔ] ‘sing’ vs.
		/loɗːo/ [ˈlɔʔɗɔ] ‘day, sun).
		\citeauthor{Onvlee1984}’s orthographic practice is to mark the vowel length
		rather than the consonant length.
		Indeed, given the distribution of
		lengthened consonants and long vowels,
		an alternate analysis would be to see vowel length as distinctive
		with consonants automatically doubled/lengthened after short vowels.}
However, \citeauthor{Onvlee1984} does not mark
lengthened consonants when the preceding vowel is high,
e.g. Kodi /ɓiɲːa/ ‹binya› ‘door’ (p. 414) vs.
/iɲa/ ‹inya› ‘mother’ (p.
107) and /rutːo/ ‹ruto› ‘blood’ (p. 439) vs.
/mutu/ ‹mutu› ‘burn’ (p. 295).

\citet{Onvlee1984} transcribes voiced plosives differently in different languages.
Languages of Sumba have only two series of voiced plosives: either
implosive and prenasalised (e.g. /ɓ/ vs.
/mb/), or implosive and plain voiced (e.g. /ɓ/ vs. /b/).
The only language with the latter system in \citet{Onvlee1984} is Anakalang.
For Anakalang \citeauthor{Onvlee1984} transcribes plain voiced stops with normal
letters and implosives with a dot beneath the letter: e.g.
‹ḅ› = /ɓ/, but ‹b› = /b/.
For other languages, \citet{Onvlee1984} transcribes implosives
with normal letters: e.g. ‹b› = /ɓ/.

To summarise, \citet{Onvlee1984} does not represent vowel or
consonant length consistently in the languages of Sumba.\footnote{
		While \citeauthor{Onvlee1984}’s transcription of implosives differs between
		different languages, the phonemic systems mean there is usually no doubt as to the
		phonetic quality specified by his symbols.}
Whenever we cite
forms from \citet{Onvlee1984} for which the underlying form is
in doubt, we replicate the transcription
and enclose it in angle brackets ‹ › to indicate that this is
a non-phonemic orthographic form.
In other cases, we have re-transcribed data from the languages
of Sumba according to their phonemic form.

